\scanpage{88}\mspage{43}
\deleted{Pour étudier l'étoile d'un}
La propriété (35) montre que les \deleted{$P_{M,i}=$}
\[
P_{M,i}=(S_M,A_{M,i},F_{M,i}\inserted{+\text{ relations d'incidence (\unsure{inclusion|induites|l'inclusion})}})\qquad(i\in\pi)\tag*{(38)}
\]
sont des structures de 2-polyèdres triangulés, sauf qu'il n'est pas clair que l'étoile de chaque sommet soit connexe
(en tant que 1-contour). Choisissons
un $e\in i=\{\pm e\}\subset\overrightarrow\pi=\det(M)^*$, on
trouve que le contour combinatoire orienté
correspondant à un sommet $s=\dot x$ est
donné par
\[
\begin{aligned}
\underbrace{\overrightarrow{\mathrm{Et}}_e(s)}_{\substack{\text{étoile}\\\text{orientée,}\\\text{définie par l'}\\\text{orientation }e\\\text{de }P_{M,i}}}
\simeq{}&\text{contour défini par l'espace}\\
&\text{affine }L_x\subset M\ \text{formé des}\\
&y\text{ tels que }x\wedge y=e,\\
&\text{ayant comme module des}\\
&\text{translations }V_x=\{u\in M\mid x\wedge u=0\}\\
&\hspace{6em}=\Lambda x,\\
&\text{avec comme ens. de}\\
&\text{sommets }L_x,\ \text{et le successeur}\\
&\text{de }y\in L_x\text{ étant }y+\deleted{\unsure{x|1}}\unsure{1|x}
\end{aligned}\tag*{(39)}
\]
% slanted
\margin{\uline{NB} \deleted{\unsure{P|Pour}} A tout $s\in S_M$, le choix de $e\in i$ définit sur \deleted{est associé \unsure{un $\Lambda$|son $\Lambda$}} $\mathrm{\acute Et}_e(s)$ une structure de \uline{torseur} \deleted{module libre de rang 1} \uline{sous $\Lambda$} (de telle façon que le \deleted{$V_s$ de $M$, $V_s=\Lambda\cdot\{x\}=\Lambda\cdot(-x)$} successeur de $y$ soit $y+1$), \deleted{(où $s=\dot x$), et $\mathrm{\acute Et}_e(s)$ est} qui est remplacée par la structure opposée si $e$ remplacé par $-e$. Donc $\mathrm{\acute Et}(s)$ est un torseur sous $\Lambda\wedge_{\mu_2}i$ ($\Lambda$ tordu par $i$)}

On trouve donc
\[
\begin{aligned}
&\text{Les étoiles des points de }P_{M,i}\text{ sont connexes}\\
&\deleted{\unsure{\text{ssi}}}\text{ ssi on a }\Lambda_0=\Lambda\quad\text{i.e. }\Lambda\text{ est un anneau absolu}\ ;\\
&\Lambda\simeq\mathbb Z\text{ ou }\Lambda\simeq\mathbb Z/n\mathbb Z.\quad\text{Dans \deleted{ce cas} \inserted{tous les cas}, les}\\
&\text{ordres des \deleted{\unsure{points}}sommets sont tous égaux à }\operatorname{card}\Lambda
\end{aligned}\tag*{(40)}
\]
% archive sheet number: 87
